% The pubviz technical report. Build: make report (from the repository root).
% Every table, figure and diagram is made with pubviz; every number in them
% comes from report/literature.json (values quoted from papers, with their
% locations) or report/data.json (computed by report/analysis.py).
\documentclass[11pt]{article}
\usepackage[letterpaper, margin=1in]{geometry}
\usepackage[T1]{fontenc}
\usepackage{newtxtext, newtxmath}
\usepackage{graphicx}
\usepackage{amsmath}
\usepackage[round, authoryear]{natbib}
\usepackage[labelfont=bf, font=small]{caption}
\usepackage{subcaption}
\usepackage[section]{placeins}
\usepackage[hidelinks]{hyperref}
% Let a page hold more floats, so a section's exhibits stay near its text.
\renewcommand\topfraction{0.9}
\renewcommand\bottomfraction{0.8}
\renewcommand\textfraction{0.07}
\renewcommand\floatpagefraction{0.75}
\setcounter{topnumber}{3}
\setcounter{bottomnumber}{2}
\setcounter{totalnumber}{5}
\usepackage{xurl}
\usepackage[diagrams]{pubviz}
% The five charts of the decision tree in docs/decision-tree.md. Input by
% examples/gallery/tree.tex (docs images) and report/report.tex.
% Needs \usepackage[diagrams]{pubviz}.
\tikzset{
  whole words/.style={execute at begin node={\hyphenpenalty=10000\exhyphenpenalty=10000}},
  % A question; an exhibit and the pubviz call that draws it.
  q/.style={pv block, text width=27mm, whole words},
  ex/.style={pv block, pv fill=pv-band, text width=31mm, whole words},
  % A pointer to another chart.
  go/.style={pv label, text=pv-muted},
  ans/.style={font=\sffamily\scriptsize, text=pv-ink, align=left, inner sep=1pt, whole words},
}
\newcommand\treecode[1]{{\ttfamily\scriptsize #1}}
\newcommand\treenote[1]{\textcolor{pv-muted}{\scriptsize[#1]}}
% \answer[label width]{question}{target}{answer}: an arrow from a question to
% the next question or exhibit, with the answer written above it.
\newcommand\answer[4][29mm]{%
  \draw[pv flow] (#2.east) -- ++(3mm,0) |-
    node[ans, text width=#1, pos=0.5, anchor=south west, xshift=1mm] {#4} (#3.west);}

\newcommand\pvtreeStart{%
\begin{tikzpicture}[pv diagram]
  \node[q] (take) {What should the reader take away?};
  \node[ex, right=34mm of take, yshift=48mm] (dia) {Diagram \treenote{1}\\\treecode{pvdiagram}};
  \node[ex, right=34mm of take, yshift=34mm] (tab) {Table \treenote{2}\\\treecode{pvtable}; chart 2};
  \node[q, right=34mm of take, yshift=-12mm] (num) {What is each number?};
  \node[go, right=34mm of take, yshift=-58mm] (model) {chart 5};
  \answer{take}{dia}{how the method works}
  \answer{take}{tab}{exact values, or many measures per method}
  \answer[32mm]{take}{num}{a shape or a comparison}
  \answer{take}{model}{what the model does with its inputs}
  \node[go, right=34mm of num, yshift=30mm] (amt) {chart 3};
  \node[go, right=34mm of num, yshift=18mm] (runs) {chart 4};
  \node[ex, right=34mm of num, yshift=5mm] (rate) {Diverging stacked bars \treenote{3}\\\treecode{pv.ratings}};
  \node[ex, right=34mm of num, yshift=-9mm] (curve) {Lines named at their ends; panels if many \treenote{4}\\\treecode{pv.curves}};
  \node[ex, right=34mm of num, yshift=-24mm] (scale) {Log axes and a fitted power law \treenote{5}\\\treecode{pv.scaling}};
  \node[ex, right=34mm of num, yshift=-38mm] (trade) {Points and their frontier \treenote{6}\\\treecode{pv.tradeoff}};
  \answer{num}{amt}{one measured amount: a count, a share, one run}
  \answer{num}{runs}{an estimate from repeated runs: seeds, people, test items}
  \answer{num}{rate}{answers on an ordered rating scale}
  \answer{num}{curve}{a series over steps or time}
  \answer{num}{scale}{a measure against model size, data or compute}
  \answer{num}{trade}{two measures that pull against each other}
\end{tikzpicture}}
\newcommand\pvtreeTables{%
\begin{tikzpicture}[pv diagram]
  \node[q] (base) {Is the claim a change from one baseline?};
  \node[ex, right=18mm of base, yshift=9mm] (delta) {Signed changes beside the values \treenote{7}\\\treecode{delta=absolute}};
  \node[q, right=18mm of base, yshift=-11mm] (grid) {Does the pattern across cells matter more than the best one?};
  \answer[12mm]{base}{delta}{yes}
  \answer[12mm]{base}{grid}{no}
  \node[ex, right=18mm of grid, yshift=8mm] (heat) {Shading by value \treenote{8}\\\treecode{heat=group}};
  \node[ex, right=18mm of grid, yshift=-8mm] (rank) {Best and second marked, with a key \treenote{9}\\\treecode{ranks=2}, \treecode{\textbackslash pvkey}};
  \answer[12mm]{grid}{heat}{yes}
  \answer[12mm]{grid}{rank}{no}
\end{tikzpicture}}
\newcommand\pvtreeAmounts{%
\begin{tikzpicture}[pv diagram]
  \node[q] (small) {Are the differences small next to the values?};
  \node[ex, right=34mm of small, yshift=30mm] (diff) {The differences themselves \treenote{10}\\\treecode{pv.deltas}};
  \node[q, right=34mm of small, yshift=-20mm] (fac) {How many factors?};
  \answer{small}{diff}{yes: bars from zero would hide them}
  \answer{small}{fac}{no}
  \node[ex, right=34mm of fac, yshift=26mm] (bars) {Bars from zero \treenote{11}\\\treecode{pv.bars}};
  \node[ex, right=34mm of fac, yshift=12mm] (ranked) {A dot per item, sorted; not a radar chart \treenote{12}\\\treecode{pv.ranked}};
  \node[ex, right=34mm of fac, yshift=-3mm] (heatmap) {Heatmap \treenote{13}\\\treecode{pv.heatmap}};
  \node[ex, right=34mm of fac, yshift=-16mm] (group) {Grouped bars from zero \treenote{11}\\\treecode{pv.bars}};
  \node[go, right=34mm of fac, yshift=-28mm] (table) {a table: chart 2};
  \answer{fac}{bars}{one, a few items}
  \answer{fac}{ranked}{one, many items: languages, classes}
  \answer{fac}{heatmap}{two, both ordered: sizes, rates, steps}
  \answer{fac}{group}{two kinds, a few of each}
  \answer{fac}{table}{two kinds, many of each}
\end{tikzpicture}}
\newcommand\pvtreeRuns{%
\begin{tikzpicture}[pv diagram]
  \node[q] (claim) {What is the claim about?};
  \node[q, right=34mm of claim, yshift=23mm] (pair) {Did each method run on the same seeds or items?};
  \node[ex, right=34mm of claim, yshift=-2mm] (level) {Point and interval, or a gradient of the mean \treenote{15}\\\treecode{pv.means}};
  \node[q, right=34mm of claim, yshift=-25mm] (n) {How many runs per group?};
  \answer{claim}{pair}{a difference between methods}
  \answer{claim}{level}{the level of each method}
  \answer{claim}{n}{how the runs spread}
  \node[ex, right=34mm of pair, yshift=9mm] (paired) {The change per run, with its interval \treenote{14}\\\treecode{pv.deltas}, \treecode{pv.mean\_ci}};
  \node[ex, right=34mm of pair, yshift=-10mm] (indep) {The difference, with an interval for two groups \treenote{14}\\\treecode{pv.deltas}};
  \answer{pair}{paired}{yes: paired}
  \answer{pair}{indep}{no: independent}
  \node[ex, right=34mm of n, yshift=9mm] (dots) {Every run as a dot, the mean marked \treenote{16}\\\treecode{pv.dots}};
  \node[ex, right=34mm of n, yshift=-10mm] (box) {Box, strip, violin or histogram, by what is judged \treenote{17}\\\treecode{pv.distributions}};
  \answer{n}{dots}{a handful, under ten}
  \answer{n}{box}{many}
\end{tikzpicture}}
\newcommand\pvtreeModel{%
\begin{tikzpicture}[pv diagram]
  \node[q] (see) {What should the reader see of the model?};
  \node[ex, right=34mm of see, yshift=44mm] (conf) {Confusion matrix, rows as shares \treenote{18}\\\treecode{pv.confusion}};
  \node[ex, right=34mm of see, yshift=28mm] (roc) {ROC; with PR when positives are rare; DET for verification \treenote{19}\\\treecode{pv.roc}, \treecode{pv.det}};
  \node[ex, right=34mm of see, yshift=12mm] (rel) {Reliability diagram with cases per bin \treenote{20}\\\treecode{pv.reliability}};
  \node[q, right=34mm of see, yshift=-20mm] (out) {What are they?};
  \node[q, right=34mm of see, yshift=-44mm] (wt) {Weights over what?};
  \node[ex, right=34mm of see, yshift=-3mm] (emb) {A map of points, classes named on it \treenote{21}\\\treecode{pv.embedding}};
  \answer{see}{conf}{which classes it mixes up}
  \answer{see}{roc}{its errors at every threshold}
  \answer{see}{rel}{whether its probabilities can be trusted}
  \answer{see}{emb}{where it places its inputs}
  \answer{see}{out}{its outputs}
  \answer{see}{wt}{what its outputs rest on}
  \node[ex, right=18mm of out, yshift=6mm] (img) {Image grid, one region enlarged \treenote{22}\\\treecode{pv.images}};
  \node[ex, right=18mm of out, yshift=-7mm] (spec) {Spectrogram in dB, pitch on top \treenote{23}\\\treecode{pv.spectrogram}};
  \answer[12mm]{out}{img}{images}
  \answer[12mm]{out}{spec}{sounds}
  \node[ex, right=18mm of wt, yshift=1mm] (att) {Weights matrix \treenote{24}\\\treecode{pv.attention}};
  \node[ex, right=18mm of wt, yshift=-11mm] (ovl) {Map over the image in grey \treenote{25}\\\treecode{pv.overlay}};
  \node[ex, right=18mm of wt, yshift=-23mm] (tok) {Shaded words \treenote{26}\\\treecode{\textbackslash pvshade}};
  \answer[12mm]{wt}{att}{another sequence}
  \answer[12mm]{wt}{ovl}{an image}
  \answer[12mm]{wt}{tok}{the words of a text}
\end{tikzpicture}}

% The Transformer (Vaswani et al., 2017), drawn with the pubviz diagram
% styles to show a dense architecture figure in the shared colors. Needs
% \usepackage[diagrams]{pubviz}. \transformer[<options>] draws it; the
% options go to pvdiagram, for example
%
%   \transformer                                         blocks filled by kind
%   \transformer[pv kinds=false, ffn/.append style={pv new}]
%                                                        by role: only a new
%                                                        feed-forward block is filled
%
% Every node is placed relative to another, so a longer label moves its
% neighbours instead of overlapping them.
\newcommand\transformer[1][]{%
\begin{pvdiagram}[
  pv kinds,
  block/.style={pv block, minimum width=20mm, minimum height=2.6em},
  thin/.style={pv block, minimum width=20mm, minimum height=1.4em, inner ysep=1.5pt},
  % One kind per sort of layer; with pv kinds each gets its own light fill.
  attention/.style={pv kind=1}, ffn/.style={pv kind=2}, norm/.style={pv kind=3},
  embedding/.style={pv kind=4}, linear/.style={pv kind=5}, softmax/.style={pv kind=6},
  #1]
  % ---- encoder, bottom to top
  \node[pv label] (in) {Inputs};
  \node[block, embedding, above=4mm of in] (iemb) {Input\\Embedding};
  \node[pv op, above=4mm of iemb] (iadd) {};
  \node[pv wave, left=1.5mm of iadd] (ipe) {};
  \node[pv label, left=0.5mm of ipe] {Positional\\Encoding};
  \node[block, attention, above=12mm of iadd] (emha) {Multi-Head\\Attention};
  \node[thin, norm, above=1mm of emha] (en1) {Add \& Norm};
  \node[block, ffn, above=8mm of en1] (eff) {Feed\\Forward};
  \node[thin, norm, above=1mm of eff] (en2) {Add \& Norm};
  % ---- decoder, bottom to top
  \node[pv label, anchor=north] (out) at ($(in.north)+(36mm,0)$) {Outputs\\(shifted right)};
  \node[block, embedding, above=4mm of out] (oemb) {Output\\Embedding};
  \node[pv op, above=4mm of oemb] (oadd) {};
  \node[pv wave, right=1.5mm of oadd] (ope) {};
  \node[pv label, right=0.5mm of ope] {Positional\\Encoding};
  \node[block, attention, above=12mm of oadd, minimum height=3.7em] (dmha1) {Masked\\Multi-Head\\Attention};
  \node[thin, norm, above=1mm of dmha1] (dn1) {Add \& Norm};
  \node[block, attention, above=12mm of dn1] (dmha2) {Multi-Head\\Attention};
  \node[thin, norm, above=1mm of dmha2] (dn2) {Add \& Norm};
  \node[block, ffn, above=8mm of dn2] (dff) {Feed\\Forward};
  \node[thin, norm, above=1mm of dff] (dn3) {Add \& Norm};
  % ---- the repeated layers: a group box each, marked N times
  \coordinate (ebottom) at ($(emha.south)+(0,-6.5mm)$);
  \coordinate (dbottom) at ($(dmha1.south)+(0,-6.5mm)$);
  \begin{scope}[on background layer]
    \node[pv group, fit=(ebottom)(emha)(en2), inner xsep=5mm, inner ysep=3mm] (eg) {};
    \node[pv group, fit=(dbottom)(dmha1)(dn3), inner xsep=5mm, inner ysep=3mm] (dg) {};
  \end{scope}
  \node[pv label, left=1.5mm of eg] {N$\times$};
  \node[pv label, right=1.5mm of dg] {N$\times$};
  % ---- output head
  \node[thin, linear, above=5mm of dg] (lin) {Linear};
  \node[thin, softmax, above=3mm of lin] (soft) {Softmax};
  \node[pv label, above=4mm of soft] (prob) {Output\\Probabilities};
  % ---- encoder edges
  \draw[pv flow] (in) -- (iemb);
  \draw[pv flow] (iemb) -- (iadd);
  \draw[pv line] (ipe) -- (iadd);
  \coordinate (ef1) at ($(emha.south)+(0,-4mm)$);    % queries, keys, values split here
  \coordinate (er1) at ($(emha.south)+(0,-6mm)$);    % the residual leaves here
  \draw[pv line] (iadd) -- (ef1);
  \draw[pv flow] (ef1) -- (emha);
  \draw[pv flow] (ef1) -| ($(emha.south)+(-6mm,0)$);
  \draw[pv flow] (ef1) -| ($(emha.south)+(6mm,0)$);
  \draw[pv flow] (er1) -| ($(en1.west)+(-3.5mm,0)$) -- (en1);
  \coordinate (er2) at ($(en1.north)!0.4!(eff.south)$);
  \draw[pv flow] (en1) -- (eff);
  \draw[pv flow] (er2) -| ($(en2.west)+(-3.5mm,0)$) -- (en2);
  % ---- decoder edges
  \draw[pv flow] (out) -- (oemb);
  \draw[pv flow] (oemb) -- (oadd);
  \draw[pv line] (ope) -- (oadd);
  \coordinate (df1) at ($(dmha1.south)+(0,-4mm)$);
  \coordinate (dr1) at ($(dmha1.south)+(0,-6mm)$);
  \draw[pv line] (oadd) -- (df1);
  \draw[pv flow] (df1) -- (dmha1);
  \draw[pv flow] (df1) -| ($(dmha1.south)+(-6mm,0)$);
  \draw[pv flow] (df1) -| ($(dmha1.south)+(6mm,0)$);
  \draw[pv flow] (dr1) -| ($(dn1.east)+(3.5mm,0)$) -- (dn1);
  % Cross-attention: keys and values from the encoder, queries from below.
  \coordinate (df2) at ($(dmha2.south)+(0,-4mm)$);
  \coordinate (dr2) at ($(dn1.north)+(0,2mm)$);
  \draw[pv flow] (dn1) -- ($(df2)+(0,-3mm)$) -| ($(dmha2.south)+(6mm,0)$);
  \draw[pv flow] (dr2) -| ($(dn2.east)+(3.5mm,0)$) -- (dn2);
  \coordinate (etop) at ($(en2.north |- eg.north)+(0,3mm)$);
  \coordinate (gap) at ($(eg.east)!0.5!(dg.west)$);
  \draw[pv flow] (en2) -- (etop) -| (gap |- df2) -| (dmha2.south);
  \draw[pv flow] (en2) -- (etop) -| (gap |- df2) -| ($(dmha2.south)+(-6mm,0)$);
  \coordinate (dr3) at ($(dn2.north)!0.4!(dff.south)$);
  \draw[pv flow] (dn2) -- (dff);
  \draw[pv flow] (dr3) -| ($(dn3.east)+(3.5mm,0)$) -- (dn3);
  \draw[pv flow] (dn3) -- (lin);
  \draw[pv flow] (lin) -- (soft);
  \draw[pv flow] (soft) -- (prob);
\end{pvdiagram}}

% A model whose middle blocks run several times with shared weights, drawn
% in the light diagram style: thin grey links, stacks of layers as small
% pictures, the repeated part in an accent region. Needs
% \usepackage[diagrams]{pubviz}. \looped[<options>] draws it; the options go
% to pvdiagram. It is drawn for the text width (a figure*).
%
% Each stage is placed relative to the one before, so changing a gap or a
% label moves the rest with it.
\newcommand\looped[1][]{%
\begin{pvdiagram}[node distance=6mm and 9mm, #1]
  % ---- the stage that runs once before the loop
  \node[pv stack] (embed) {};
  \node[pv title, above=2.5mm of embed] {Embedding\\\pvnote{runs once}};
  % ---- the shared blocks, one row per step, in an accent region
  \foreach \step/\below in {1/embed, 2/b1, 3/b2} {
    \ifnum\step=1
      \node[pv stack, pv emph, right=27mm of embed] (a\step) {};
    \else
      \node[pv stack, pv emph, below=11mm of a\the\numexpr\step-1\relax] (a\step) {};
    \fi
    \node[pv stack, pv emph, right=7mm of a\step] (b\step) {};
    \draw[pv link, pv emph] (a\step) -- (b\step);
    \begin{scope}[on background layer]
      \node[pv region=pv-fill-3, fit=(a\step)(b\step), inner xsep=3mm, inner ysep=2.5mm] (r\step) {};
    \end{scope}
  }
  \node[pv title, above=2.5mm of r1] {Shared blocks\\\pvnote{\pvrepeat\ run three times}};
  % ---- the state each step hands to the next
  \foreach \from/\to in {1/2, 2/3} {
    \draw[pv link, pv emph] (r\from.south) -- node[pv note, right, text=pv-strong] {state} (r\to.north);
  }
  % ---- the input reaches every step along one branching line
  % Each branch is drawn whole from the start, so its turn down is rounded.
  \coordinate (bus) at ($(embed.east)+(7mm,0)$);
  \draw[pv link] (embed.east) -- node[pv note, above, pos=0.78] {Step 1} (r1.west);
  \draw[pv link] (embed.east) -- (bus) |- node[pv note, above, pos=0.8] {Step 2} (r2.west);
  \draw[pv link] (embed.east) -- (bus) |- node[pv note, above, pos=0.8] {Step 3} (r3.west);
  % ---- the stage that runs once after the loop
  \node[pv stack, right=24mm of r3] (head) {};
  \node[pv title, above=2.5mm of head] {Output head\\\pvnote{runs once}};
  \draw[pv link] (r3.east) -- (head);
\end{pvdiagram}}

% A speech recognizer as a pipeline, left to right, in the light diagram
% style: a small picture for each kind of data, stacks of layers for the
% networks, a title and a grey note above each stage on one line, the new
% encoder in the accent, and the path used only in training dashed. Needs
% \usepackage[diagrams]{pubviz}. \asr[<options>] draws it; the options go to
% pvdiagram. It is drawn for the text width (a figure*).
\newcommand\asr[1][]{%
\begin{pvdiagram}[node distance=6mm and 12mm,
  stage/.style={minimum width=14mm, minimum height=11mm, inner sep=0pt},
  #1]
  % ---- what goes in: a waveform, then the features computed from it
  \node[stage] (wave) {};
  \draw[pv-line, line width=0.5pt, shift={(wave.center)}] plot[domain=-1:1, samples=90, smooth]
    ({7mm*(\x)}, {4mm*sin(1100*(\x))*exp(-3*(\x)^2)});
  \node[stage, right=of wave] (feat) {};
  \foreach \c [count=\i from 0] in {20,45,70,40,15,35,60,80,50,25,10,30,55,65,45,20,
                                    15,25,40,30,20,35,50,60,35,15,5,10,20,25,15,10} {
    \pgfmathtruncatemacro\col{mod(\i,8)}
    \pgfmathtruncatemacro\row{div(\i,8)}
    \fill[pv-line!\c!white] ($(feat.south west)+(\col*1.75mm, 0.5mm+\row*2.5mm)$) rectangle ++(1.6mm, 2.35mm);
  }
  % ---- the networks: the new encoder in the accent
  \node[pv stack, pv emph, right=of feat] (enc) {};
  \node[pv stack, right=of enc] (dec) {};
  % ---- what comes out
  \node[stage, right=of dec] (text) {\itshape ``the cat sat''};
  % ---- one line of titles above every stage
  \coordinate (titles) at ($(enc.north)+(0,2.5mm)$);
  \foreach \at/\title in {wave/{Speech\\\pvnote{16 kHz}}, feat/{Features\\\pvnote{80 mel bands}},
                          enc/{\textbf{Encoder}\\\pvnote{new; 12 layers}}, dec/{Decoder\\\pvnote{6 layers}},
                          text/{Text\\\pvnote{tokens}}}
    \node[pv title, anchor=south] at (\at |- titles) {\title};
  \draw[pv link] (wave) -- (feat);
  \draw[pv link] (feat) -- (enc);
  \draw[pv link, pv emph] (enc) -- node[pv note, above, text=pv-strong] {states} (dec);
  \draw[pv link] (dec) -- (text);
  % ---- a path taken only in training, dashed, below the main row
  \draw[pv link, dashed] (text.south) -- ++(0,-6mm) -| node[pv note, below, pos=0.25]
    {CTC loss, training only} (enc.south);
\end{pvdiagram}}

% Retrieval-augmented generation in the light diagram style: a question, a
% retriever that searches an index, the passages it finds, and a reader
% that writes the answer from them and the question. The retrieval stage
% sits in a neutral region, and the reader, the part this example takes as
% new, is in the accent. Needs \usepackage[diagrams]{pubviz}; the index
% is a cylinder from the TikZ library loaded here. \retrieval[<options>]
% draws it; the options go to pvdiagram. It is drawn for the text width (a
% figure*).
\usetikzlibrary{shapes.geometric}
\newcommand\retrieval[1][]{%
\begin{pvdiagram}[node distance=6mm and 12mm,
  stage/.style={minimum width=14mm, minimum height=11mm, inner sep=2pt, align=center},
  doc/.style={draw=pv-line, fill=white, minimum width=6mm, minimum height=7.5mm, inner sep=0pt,
    path picture={\foreach \y in {0.25,0.45,0.65}
      \draw[pv-muted, line width=0.4pt] ($(path picture bounding box.south west)!\y!(path picture bounding box.north west)+(1.2mm,0)$)
        -- ++(3.6mm,0);}},
  #1]
  \node[stage] (q) {Who wrote\\``Middlemarch''?};
  % ---- retrieval: a retriever and the index it searches
  \node[pv stack, pv layers=4, right=of q] (ret) {};
  \node[cylinder, shape border rotate=90, aspect=0.25, draw=pv-line, fill=white,
    minimum width=11mm, minimum height=9mm, below=8mm of ret] (index) {};
  \node[pv note, anchor=west] at ($(index.east)+(1.5mm,0)$) {21M\\passages};
  \draw[pv link] ([xshift=-1.5mm]ret.south) -- node[pv note, left] {search} ([xshift=-1.5mm]index.north);
  \draw[pv link] ([xshift=1.5mm]index.north) -- ([xshift=1.5mm]ret.south);
  % ---- the passages found, then the reader and its answer
  \node[stage, right=of ret] (docs) {};
  \node[doc] (d1) at ($(docs.center)+(-2mm,1.5mm)$) {};
  \node[doc] (d2) at (docs.center) {};
  \node[doc] (d3) at ($(docs.center)+(2mm,-1.5mm)$) {};
  \node[pv stack, pv emph, right=22mm of docs] (read) {};
  \node[stage, right=of read] (a) {George Eliot};
  \begin{scope}[on background layer]
    \node[pv region, fit=(ret)(index), inner xsep=7mm, inner ysep=3mm] {};
  \end{scope}
  % ---- one line of titles above every stage, aligned at the top
  \coordinate (titles) at ($(ret.north)+(0,11mm)$);
  \foreach \at/\title in {q/{Question}, ret/{Retriever\\\pvnote{dual encoder}}, docs/{Passages\\\pvnote{top 5}},
                          read/{\textbf{Reader}\\\pvnote{new}}, a/{Answer}}
    \node[pv title, anchor=north] at (\at |- titles) {\title};
  \draw[pv link] (q) -- (ret);
  \draw[pv link] (ret) -- (docs);
  \draw[pv link] (docs) -- node[pv note, below] {and the question} (read);
  \draw[pv link, pv emph] (read) -- (a);
\end{pvdiagram}}


\newcommand\build{../examples/build/report}
\newcommand\fig[2][\linewidth]{\IfFileExists{\build/#2.pdf}
  {\includegraphics[width=#1]{\build/#2}}{\fbox{figure \detokenize{#2}}}}
\newcommand\gen[1]{\IfFileExists{\build/#1.tex}{\input{\build/#1.tex}}{\fbox{table #1}}}
\newcommand\code[1]{\texttt{#1}}

\title{\textbf{pubviz}: Choosing and Drawing the Tables, Figures\\and Diagrams of a Research Paper}
\author{Jhen-Ke Lin\\National Yang Ming Chiao Tung University\\\texttt{jacob.cs14@nycu.edu.tw}}
\date{}

\begin{document}
\maketitle

\begin{abstract}
Most conferences and journals provide a paper template. These templates set
the page layout but say little about how results should be shown, so each
paper decides on its own which exhibit to use and how to mark, color and size
it. Many studies have measured how readers take in charts, tables and
diagrams, but their findings reach the tools researchers use slowly. We
present pubviz, an open library that brings these findings into a paper in
two steps. A decision tree chooses the exhibit from what the numbers are and
what the reader should take away, for a model's results as well as for what
it does with its inputs, and ties each answer to the study behind it. A LaTeX package and a Python library then draw the chosen table, figure
or diagram with defaults based on these findings: compared values on one
scale, the uncertainty of an estimate shown over the runs themselves,
colormaps ordered in lightness, fills that survive a grey print, and figures
sized from the paper's own layout. Every exhibit takes its colors from one
accent and keeps its meaning in grey and without color, and its words follow
the paper's language. Where we found no study, as for how a table should mark
its best value, we say so and mark the choice as our judgment. Every exhibit
in this report is made with pubviz.
\end{abstract}

\section{Introduction}\label{sec:intro}

Most conferences and journals provide a paper template. The template fixes
the page: margins, columns, fonts and the style of captions. It says almost
nothing about the results on that page: whether a comparison should be a
table or a figure, how a table points out its best number, how a figure sets
the proposed method apart, or what its colors become on a grey printer. Each
paper settles these questions on its own.

The questions are not new to research. Experiments have measured how readers
judge positions, lengths and angles
\citep{cleveland1984graphical,heer2010crowdsourcing}, how they read error bars
and intervals \citep{belia2005researchers,correll2014error,hofman2020visualizing},
and how colormaps and palettes help or mislead them
\citep{liu2018rainbow,szafir2018modeling,schloss2019mapping}. New results
reach practical tools slowly, as \citet[Sec.~1]{moritz2019draco} note for
visualization tools in general, and a researcher writing a paper meets them,
if at all, as scattered rules of thumb.

Papers in machine learning also show what a model does, not only how well it
scores: the classes it confuses, the images or sounds it produces, the
weights it places on its inputs. Guidance on charts mostly covers plots of
values, yet in two large counts of computer science papers the largest share
of figures fell outside the plot and diagram types that were counted
\citep[Table~1]{siegel2016figureseer}\citep[Sec.~3]{hsu2021scicap}. For these
exhibits, studies with readers are fewer still, and much of what is known
comes from methodological arguments and from studies in other fields.

We present pubviz, an open library that brings these findings into a paper
in two steps (Figure~\ref{fig:overview}). The first step chooses the exhibit.
A decision tree asks what the reader should take away, what each number is,
how many factors it varies over and how many values there are; each answer
leads to the next question or to a table, figure or diagram of a named kind.
Every answer is tied to the study it rests on, or marked as our judgment
where we found none. The second step draws the chosen exhibit. A LaTeX
package marks result tables and draws method diagrams inside the paper, and
a Python library draws the figures. Both take their defaults from the same
findings and their colors from one accent color, and the LaTeX side passes
the page's measurements to the figures through a layout file.

\begin{figure}[tb]
  \centering
  \begin{pvdiagram}[node distance=5mm and 9mm]
    \node[pv label] (claim) {A claim\\and its data};
    \node[pv block, right=of claim, minimum width=31mm] (tree) {Choose the exhibit\\decision tree, \S\ref{sec:tree}};
    \node[pv block, right=12mm of tree, yshift=8mm, minimum width=35mm] (tex) {Tables and diagrams\\LaTeX package};
    \node[pv block, right=12mm of tree, yshift=-8mm, minimum width=35mm] (py) {Figures\\Python library};
    \node[pv label, right=12mm of tex, yshift=-8mm] (page) {The printed\\page};
    \begin{scope}[on background layer]
      \node[pv group, fit=(tree)(tex)(py), inner sep=3mm,
        label={[pv note]above:pubviz}] {};
    \end{scope}
    \draw[pv flow] (claim) -- (tree);
    \draw[pv flow] (tree.east) -- ++(4mm,0) |- (tex.west);
    \draw[pv flow] (tree.east) -- ++(4mm,0) |- (py.west);
    \draw[pv flow] (tex.east) -- ++(4mm,0) |- (page.west);
    \draw[pv flow] (py.east) -- ++(4mm,0) |- (page.west);
    \draw[pv flow] (tex.south) -- node[pv note, left, align=right] {layout file: widths,\\fonts, colors, words} (py.north);
  \end{pvdiagram}
  \caption{\textbf{pubviz first chooses the table, figure or diagram,
    then draws it.} The decision tree (Section~\ref{sec:tree}) picks a
    table, figure or diagram; the LaTeX package and the Python library draw
    it (Sections~\ref{sec:design} and~\ref{sec:catalog}). Each compile of the
    paper writes a layout file that sizes and colors the figures.}
  \label{fig:overview}
\end{figure}

This report makes four contributions:
\begin{itemize}
  \item A decision tree for the exhibits of a research paper, from result
    tables to the exhibits that show what a model does, in which each
    answer points to the place in a study that supports it, or is marked as
    our judgment (Section~\ref{sec:tree}).
  \item A LaTeX package that marks result tables from their numbers when the
    paper compiles, with ranks, changes against a baseline, shading and a key
    that explains each mark; and a Python library that sizes and colors
    figures from a layout file the paper itself writes
    (Sections~\ref{sec:design} and~\ref{sec:catalog}).
  \item Defaults for color, uncertainty and print based on these findings,
    with three checks of our own for color choices the studies leave open
    (Sections~\ref{sec:design} and~\ref{sec:checks}).
  \item An account of what remains untested, above all how a table should
    mark its best value (Section~\ref{sec:limits}).
\end{itemize}
Every table, figure and diagram in this report is made with pubviz, and
every number in them is computed from data files in its repository.

\section{Related work}\label{sec:related}

\paragraph{Choosing an exhibit automatically.} Automatic presentation began
with APT \citep{mackinlay1986apt}, which picks an encoding from the types of
the data and ranks the candidates by the accuracy ordering of
\citet{cleveland1984graphical}. APT already designs differently for a color
and a monochrome medium \citep[Sec.~9]{mackinlay1986apt}, but its ranking for
ordinal and nominal data is, in its own words, conjectural and not
empirically verified \citep[p.~124]{mackinlay1986apt}. Show Me brings
automatic presentation into Tableau with rules taken from design books and
leaves the user's task out on purpose \citep[Sec.~1]{mackinlay2007showme}.
Voyager and its query engine CompassQL rank encodings with hand-tuned weights
after Cleveland and Mackinlay
\citep{wongsuphasawat2016voyager,wongsuphasawat2016compassql}. Draco is the
first to take a task as input and to learn its weights from perception
experiments; its authors limit the gain to the design space those
experiments covered \citep[Sec.~6.3]{moritz2019draco}. VizML learns from the
charts that Plotly users made \citep{hu2019vizml}. All five serve analysis on
a screen. None handles result tables, the uncertainty of an estimate, a grey
print or the size of a venue's column, and none asks what claim the exhibit
makes.

\paragraph{Guides for choosing a chart.} Several guides organize the choice
for people rather than programs. From Data to Viz is a decision tree by the
format of the data and cites no studies \citep{holtz2018datatoviz}.
\citet{wilke2019fundamentals} organizes a book by the kind of quantity shown
and cites studies mainly for uncertainty and color vision deficiency;
\citet{rougier2014tenrules} and \citet{midway2020principles} give ten rules
each, resting mostly on experience. The tree in Section~\ref{sec:tree} differs
in two ways: it branches on the claim as well as on the data, and it covers
tables, which these guides mostly leave aside; Wilke gives rules for the
layout of a table but none for marking its best value
\citep[Sec.~22.3]{wilke2019fundamentals}.

\paragraph{What papers show and report.} Counts of figures in computer
science papers find x--y plots the most common single type, about a fifth of
the figures in two large samples, with bar charts and scatterplots far behind
\citep[Table~1]{siegel2016figureseer}\citep[Sec.~3.1]{hsu2021scicap}; in the
ACL Anthology, tables outnumber figures \citep[Sec.~4.2]{karishma2023aclfig}.
Surveys of reporting practice find gaps that an exhibit and its caption can
close. Almost half of the natural language generation papers in one survey
did not give the number of people in a human evaluation
\citep[p.~357]{vanderlee2019best}, and most evaluations in another did not
show the prompt their evaluators saw \citep[Sec.~3.2.2]{howcroft2020twenty}. In
speech synthesis, most papers used fewer than 20 listeners
\citep[Table~1]{wester2015are} and most did not state the labels of the
rating scale \citep[Table~3]{kirkland2023stuck}. Few papers on generated text
show failures among their examples \citep[Table~1]{gehrmann2023repairing}, and
precision--recall curves were rare even in studies of unbalanced data
\citep[Table~4]{saito2015precision}.

\paragraph{Tools for paper exhibits.} Plotting libraries such as matplotlib
\citep{hunter2007matplotlib} leave the size of a figure to the user.
SciencePlots and tueplots \citep{garrett_scienceplots,kramer_tueplots} ship
styles with the widths and fonts of named venues written into their code;
matplotlib-latex-bridge compiles a sample document of a class to measure
them \citep{vannucci_mlb}; and pgfplots draws inside the document, so its
figures take the document's width and fonts by construction
\citep{feuersaenger2026pgfplots}. The closest precedent for the layout file of
pubviz is a worked example in hybrid-latex, in which LaTeX writes a
figure's width and height to a JSON file that Python reads
\citep{brewin2018hybridlatex}. For tables, booktabs, siunitx and colortbl set
rules, align numbers and fill cells
\citep{fear2020booktabs,wright2026siunitx,carlisle2026colortbl}, but the
author decides which cell to mark. pandas can bold or shade the best cell
when it writes a table out as LaTeX \citep{pandas_styler_to_latex}, outside
the paper and without a key, and xcolor can turn every color of a document to
grey to preview a black-and-white print \citep{kern2024xcolor}. Among the
items we read, we found no tool that ranks and marks a table when the paper
compiles and explains the marks in a key, none that sizes figures from a
layout written by the paper's own compilation, and none that shares one color
theme across tables, figures and diagrams. These are statements about the
items we read, not a proof that nothing else exists.

\section{Choosing an exhibit}\label{sec:tree}

The decision tree asks its questions in five flowcharts
(Figures~\ref{fig:tree-start}--\ref{fig:tree-model}); the last is for what
a model does with its inputs. A white box is a
question and the words on an arrow are one answer to it; a grey box is an
exhibit, with the pubviz call that draws it and the number of the note in
the repository's \code{docs/decision-tree.md} that gives its evidence. This
section summarizes that evidence. We keep apart three kinds of support:
studies with readers, cited with the place where the finding is; our own
computations, reported in Section~\ref{sec:checks}; and our judgment, where
we found neither. Table~\ref{tab:status} sorts the answers by that support.

\subsection{Where to start}

\begin{figure}[tb]
  \centering
  \pvtreeStart
  \caption{\textbf{The first questions separate a method from its results,
    and lookup from shape.} Chart 1 of the decision tree. Charts 2 to 5
    continue the branches marked there.}
  \label{fig:tree-start}
\end{figure}

A diagram is for how the method works. It helps when it puts information
used together in one place, so a reader finds it by location instead of by
label, and an arbitrary diagram earns none of this
\citep[pp.~98--99]{larkin1987diagram}. Arrows read as process: with arrows,
people described more of how a device works and less of its parts, and
people who drew a device after reading about its function added arrows far
more often \citep[pp.~584--588]{heiser2006arrows}. Of five layout properties
of a node-link diagram, only edge crossings affected both errors and
response time \citep[Sec.~3.1]{purchase1997aesthetic}.

For results, the first question is whether readers will look up exact values
or need the shape. On the same data and tasks, a table was among the most
accurate and fastest forms for looking up a single value and the least
accurate for judging whether two measures are correlated
\citep[Table~1]{saket2019task}. For spotting main effects and interactions in
two-by-two results, line graphs were faster than tables but not more accurate
\citep[pp.~7, 11]{braithwaite2013flexibility}. The older argument between
Ehrenberg, who held that graphs are of little use for quantitative detail
\citep[pp.~281--286]{ehrenberg1977rudiments}, and Gelman et al., who found
that all but 4 of the 64 tables in one journal issue could have been graphs
\citep[Fig.~8]{gelman2002practice}, rests on worked examples rather than
readers. We read the experiments as a
division of labor: a table when exact values or many measures matter, a
figure when the shape does. Past about six methods or six groups we send the
result to a table; that limit is our judgment.

A series over steps or time goes on one axis with each line named at
its end. Readers judge values most accurately as positions on a common scale
\citep[p.~536]{cleveland1984graphical}, and they take lines as trends: in one
study they described all 36 line graphs as trends and 28 of 33 bar graphs as
comparisons between separate values \citep[Table~1]{zacks1999bars}. Overlaid
lines were fastest when readers compared series at the same point of the
x axis, and side-by-side panels when the compared points lay far apart, with
no difference in correctness \citep[Sec.~7, Fig.~11]{javed2010graphical}. In
an informal follow-up with four people, the two layouts took about the same
time with two to four series and pulled apart from about eight
\citep[Fig.~12]{javed2010graphical}; we switch to panels there, a judgment
resting on that small test.

Two measures that pull against each other, such as quality and cost, become
one point per method with a line through the methods that no other method
beats on both. We found no study of displays for such trade-offs; the nearest
result is that scatterplots helped readers judge correlation in small
datasets \citep[Table~1]{saket2019task}. Two scales in one plot area misled:
when two scales of the same variable shared a plot, readers were least
accurate and slowest, and all but one ranked that chart lowest
\citep[Secs.~5--6]{isenberg2011study}. Readers of dual-axis line charts
described the two series as correlated far more often than readers of
connected scatterplots \citep[Sec.~3.3]{haroz2016connected}, although the
crossings of the two lines mean nothing unless the axes share units and
scale \citep[Sec.~2.2]{haroz2016connected}. pubviz therefore keeps one axis
per measure; no study has tested twin y axes for two different measures.

A measure against model size, data or compute goes on logarithmic axes with
a fitted power law. Log axes carry a cost: ecologists read log--log graphs
correctly far less often than linear ones \citep[abstract]{menge2018logarithmic},
and a general sample misjudged growth on a log axis
\citep[Table~7]{romano2020scale}. Readers who fit a trend by eye make it too
steep \citep[Sec.~3.2.2]{ciccione2021can}, so the fitted line is drawn rather
than left to the eye. Which points a fit used is the other risk: scaling-law
papers extend their fits far past the data with little or no change of style
\citep[Fig.~15]{kaplan2020scaling}\citep[Figs.~1--3]{hoffmann2022training}.
pubviz draws the fit solid over the points it used and dashed beyond, and
draws held-out points hollow, much as the GPT-4 report gives its held-out
point a color of its own \citep[Fig.~1]{openai2023gpt4}; no study has tested
these marks.

Answers on a rating scale, such as a quality score or a preference between
two systems, are ordered categories, and the tree draws how they spread
rather than their mean alone. Analyzing such answers as numbers produced
false alarms, misses and reversed orders in simulation
\citep[Figs.~2--4]{liddell2018analyzing}, yet all 68 psychology articles that
Liddell and Kruschke checked did so \citep[pp.~4--5]{liddell2018analyzing}.
pubviz draws diverging stacked bars, the layout
\citet[p.~6]{heiberger2014design} argue for because the shares on each side
of the scale then start from one line. No reader study has compared it with
other layouts; the nearest experiment, on stock data, found it no faster than
ordinary stacked bars, though rated easier \citep[Secs.~5.1, 5.3]{indratmo2018efficacy}.
When a paper reports a mean opinion score, ITU-T P.800.2 asks for the number
of votes and the standard deviation beside it and warns against comparing
scores across tests \citep[cls.~8, 10]{itu2016p8002}: one system's score fell
by up to 1.28 points when it was retested among a narrower range of systems,
in which it ranked last \citep[Sec.~4.4]{cooper2023investigating}. In win / tie / loss judgments the
ties are set aside, a choice we found no study of.

\subsection{Tables}

\begin{figure}[tb]
  \centering
  \pvtreeTables
  \caption{\textbf{A table's marks follow its claim.} Chart 2 of the
    decision tree.}
  \label{fig:tree-tables}
\end{figure}

A table whose claim is a change from a baseline prints the signed change
beside each value, so the reader need not subtract; this follows the argument
of \citet[Sec.~2]{gleicher2011visual} for drawing a relationship directly,
which they argue rather than measure. A table whose point is a pattern across
its cells shades each cell by its value. On a white page, darker is read as
more, and a mapping that runs the other way costs time even with a legend in
view \citep[Sec.~3]{schloss2019mapping}. pubviz shades better values darker
in every column, so in a column where lower is better the darker cells hold
the smaller numbers; the key states the direction, and whether that is
enough is untested.

Otherwise the table marks its best and second values and explains the marks
in a key. No study we found tested marks in table cells. On text, bold, a
yellow background and red text helped readers find highlighted words about
equally well, and bold with a background is a pair the authors recommend
because the two barely interfere \citep[Fig.~3, Secs.~6.3, 8.1]{strobelt2016guidelines};
in a word search, bold was found faster than red text
\citep[Tables~2--3]{doi2023highlighting}. Show Me leaves its color-filled
tables unranked because a fill can make the text hard to read
\citep[Sec.~5.2]{mackinlay2007showme}; the pubviz fills keep black text
well above that risk (Section~\ref{sec:checks}). The key is our inference:
unexplained marks are common and misread, as 22 of 46 visualization papers
did not say what their error bars were \citep[Sec.~2.3]{correll2014error} and
researchers misread both kinds of error bar \citep[p.~393]{belia2005researchers}.

\subsection{One measured amount per item}

\begin{figure}[tb]
  \centering
  \pvtreeAmounts
  \caption{\textbf{A number measured once is a bar from zero, unless its
    differences are what matters or the items are many.} Chart 3 of the
    decision tree, for counts, shares and scores from a single run.}
  \label{fig:tree-amounts}
\end{figure}

A number measured once, such as a count, a share or a score from a single run
on a fixed test set, is drawn as a bar from zero. Bars are designed for
categories \citep[p.~1]{weissgerber2015beyond}, and their length is their
value, so the axis must start at zero: cutting a bar axis, stretching a line
chart or drawing values as bubble areas each made readers rate the same
difference as larger (Figure~\ref{fig:pandey}), and a broken-axis mark did
not remove the effect \citep[pp.~6, 8]{correll2020truncating}. When the
differences are small next to the values, as between scores near 90, bars
from zero hide them; the way out is to draw the differences themselves, the
third strategy of \citet[Sec.~2]{gleicher2011visual}. When both factors are
ordered, as model size by data size, the grid reads as a surface and becomes
a heatmap whose colormap keeps its lightness in order
(Section~\ref{sec:checks}).

When the one factor has many items, such as the languages of a multilingual
test or the classes of a detector, bars crowd and papers often turn to radar
charts. Readers were slower on radial versions of common graphs than on
linear ones \citep[abstract]{goldberg2011eyetracking} and recalled positions
in radial grids less accurately than in rectangular ones
\citep[Sec.~4.2]{diehl2010radial}; among radial designs the radar chart did
worst \citep[abstract]{albo2016radar}, and a review of glyph studies found
linear designs ahead in most comparisons \citep[Sec.~4.6.4]{fuchs2017glyphs}.
pubviz draws a dot chart instead, the form \citet{cleveland1984methods}
proposed to replace bar charts, with the rows sorted by the paper's own
method; the sorting and the marker per system are our judgment.

\begin{figure}[tb]
  \centering
  \fig{pandey}
  \caption{\textbf{Each distortion made the same difference look larger.}
    Change in the mean rating of how much better one quantity was than
    another (five-point scale), distorted chart minus control, from
    \citet[Tables~1--2]{pandey2015deceptive}. Lines: approximate 95\%
    intervals of the change, which we derived from the two reported intervals
    assuming independent groups \citep{cumming2005inference}. Each distortion
    had its own control, so the rows are not meant to be compared with each
    other. Worse: the chart overstated the difference.}
  \label{fig:pandey}
\end{figure}

\subsection{An estimate from repeated runs}

\begin{figure}[tb]
  \centering
  \pvtreeRuns
  \caption{\textbf{For an estimate from repeated runs, the claim decides
    between the difference, the level and the spread.} Chart 4 of the
    decision tree.}
  \label{fig:tree-runs}
\end{figure}

Many results in machine learning are means over seeds, and for these the
tree asks what the claim is about. When it is a difference between methods, the figure
draws that difference with its own interval. For repeated measures, the
intervals of the separate means say nothing about the difference
\citep[pp.~176--177]{cumming2005inference}, and researchers asked to place two
means ``just significantly different'' erred with both interval and
standard-error bars (Figure~\ref{fig:belia}). When the methods ran on the
same seeds, the interval is that of the per-seed changes, which
\code{pv.mean\_ci} computes; otherwise the two margins combine as for
independent groups, by the formula in the tree's note. Drawing the difference
beside the data is argued for by \citet{ho2019moving}, but how readers take a
drawn difference has not been tested \citep{hofman2020visualizing}.

\begin{figure}[tb]
  \centering
  \fig{belia}
  \caption{\textbf{Researchers put interval bars too far apart and
    standard-error bars too close.} The $p$ implied by where respondents
    placed the second of two means so that they looked just significantly
    different, by discipline; the correct placement gives $p=.05$. Values
    from \citet[p.~393]{belia2005researchers}.}
  \label{fig:belia}
\end{figure}

When the claim is the level of each method, the figure draws a point with
its interval, never a bar: with bars, readers judged values inside the bar
likelier than values above it \citep[Secs.~4.2, 4.4]{correll2014error}. An
interval alone still misleads. Shown 95\% confidence intervals, crowd workers
put the chance that a treated case beats a control at 86\% against a true 57\%
\citep[p.~5]{hofman2020visualizing}; data scientists shown standard errors put
it at 76\%, and at 61\% when every data point was drawn too, against a true
59\% \citep[Fig.~3]{zhang2023illusion}. pubviz therefore draws the runs
faintly under the estimate by default. It can also draw the interval as a
gradient, which among researchers halved the drop in confidence around
$p=.05$, as a violin did; readers took the violin of a mean for the violin of
the data \citep[Fig.~4]{helske2021can}, so pubviz offers only the gradient.

When the claim is how the runs spread, the number of runs decides. With a
handful, every run is drawn as a dot: summary statistics are only meaningful
when there are enough data to summarize \citep[pp.~5--6]{weissgerber2015beyond},
and with five points the shape of a distribution cannot be told
\citep[Fig.~7]{weissgerber2019reveal}. With many, no display is best for
everything. With 200 points per chart, box plots were best for the median and
for comparing two groups and worst for the range, and most readers took the
box for a cluster; strip plots were best for the range and violins for
clusters \citep[Sec.~4]{lotteraner2026investigating}. Strip plots also gave
the best estimates of spread \citep[Table~1]{newburger2023fitting}. The
cut-off between the two, about ten runs, is our rule of thumb;
\citet{weissgerber2019reveal} find no clear guideline for it.

\subsection{What a model does}

\begin{figure}[tb]
  \centering
  \pvtreeModel
  \caption{\textbf{Exhibits of a model at work show its errors, its outputs
    or the weights its outputs rest on.} Chart 5 of the decision tree.}
  \label{fig:tree-model}
\end{figure}

A confusion matrix shows which classes a classifier mixes up, and its
readers are easily confused themselves. Lay readers did better when the
cells used the task's own words instead of generic terms
\citep[Table~2]{shen2020designing}; in a think-aloud study, users swapped
missed cases and false alarms when rows and columns looked alike, and found
counts easier than rates \citep[Secs.~7.4--7.5]{beauxisaussalet2019supporting}.
pubviz names the axes as the true and the predicted class, shows each row as
shares of its class so that rare and common classes compare, an argument
made without a study \citep[Sec.~4.1]{gortler2022neo}, and prints the counts
beside the shares when there are few classes. No study has tested counts
against shares, the treatment of the diagonal or the color scale
\citep[Sec.~7.1]{ren2017squares}.

Curves over every threshold rest on methodological arguments rather than on
readers. An ROC curve does not change with class balance and a
precision--recall curve does \citep[Sec.~4.2]{fawcett2006introduction}, so on
unbalanced data an ROC curve can look nearly perfect while the
precision--recall curve shows much room to improve
\citep[Fig.~1]{davis2006relationship}. pubviz draws the precision--recall curve
in steps, since straight lines between its points overstate the area
\citep[Sec.~4]{davis2006relationship}, with the share of positives as its
baseline \citep[p.~7]{saito2015precision}, and draws DET curves on
normal-deviate scales for verification, where good systems would otherwise
bunch in one corner \citep{martin1997det}. A reliability diagram hides how
many cases each bin holds \citep[Sec.~2]{guo2017calibration}, and its look
changes with the number of bins \citep[Fig.~2]{dimitriadis2021stable}; pubviz
places each point at its bin's mean prediction
\citep[Sec.~2]{brocker2007increasing} and draws the counts underneath.

Outputs that are images are compared side by side, which is slow and misses
small differences: in a flicker task, viewers needed many alternations to see
a change between two photographs until they were told where it was
\citep[p.~370]{rensink1997see}. In interactive viewers, separate windows shown
together beat zooming once a pattern had more than a few parts
\citep[Table~I]{plumlee2006zooming}, and magnified insets made comparisons more
accurate than outlined boxes \citep[Sec.~6.1]{lekschas2020pattern}. pubviz
therefore enlarges one region under every image and can add each output's
difference from the reference, as a difference image helped radiologists find
changes \citep[abstract]{difazio1997digital}. Which examples to show matters as
much as how: random examples give the least biased picture
\citep[Sec.~4.2]{vanmiltenburg2021underreporting}. For sound we found no study
of spectrogram displays. pubviz draws levels in decibels down to a fixed
floor, from white for quiet to near black for loud in even lightness steps
(Section~\ref{sec:checks}), so that darker means more as readers expect on a
white page \citep[Sec.~3]{schloss2019mapping}. In another field, a rainbow map
led medical students to find fewer of the regions they looked for even
though they preferred it \citep[Sec.~6.5]{borkin2011evaluation}.

Exhibits of what a model's outputs rest on show weights, and the weights
claim less than they seem to. Very different attention weights gave the same
output \citep[Sec.~4.2]{jain2019attention}, raw attention in the upper layers of
a Transformer is close to uniform \citep[Sec.~2]{abnar2020quantifying}, and some
saliency methods give nearly the same map for a model whose weights were
randomized \citep[Sec.~3]{adebayo2018sanity}. Readers also trust them too much:
on one dataset they found Grad-CAM explanations of wrong predictions nearly
as convincing as those of right ones \citep[Table~1]{kim2022hive}. pubviz draws
attention without numbers, lays a saliency map over the image in grey with
its opacity rising with the weight, and gives a signed map two hues, since a
map of absolute values hides where the sign flips
\citep[Fig.~4]{adebayo2018sanity}. Over text it shades the words in the
paper's own font; there readers judged longer words as more important at the
same shade \citep[Sec.~5.1]{schuff2022human}, a bias a caption can only warn
of. Maps of embeddings mislead differently. In t-SNE maps the sizes of
clusters and the distances between them may mean nothing
\citep[Secs.~2--3]{wattenberg2016tsne}, and readers found clusters well but
judged which was nearest or densest poorly
\citep[Sec.~6.3.1]{xia2022revisiting}. pubviz names each class on the map,
hides the axis values, and folds the classes beyond the colors readers can
tell apart into one grey group; in scatterplots, accuracy fell as categories
were added, above all from six on \citep[Secs.~4.1, 5.2]{tseng2023categorical}.

\begin{table}[tb]
  \centering\small
  \caption{\textbf{The choices about what to draw rest on studies; the
    choices about how to mark it rest on judgment.} The answers of the
    decision tree and the design choices of pubviz, by what supports them.}
  \label{tab:status}
  \begin{pvtable}[ranks=0, stretch=1.05]{ll}
    \toprule
    Choice & Basis \\
    \midrule
    \pvsection{Supported by studies with readers} \\
    Table for lookup, figure for shape & \citet{saket2019task}; \citet{braithwaite2013flexibility} \\
    Compared values on one positional scale & \citet{cleveland1984graphical}; \citet{heer2010crowdsourcing} \\
    Bars from zero; differences drawn when small & \citet{pandey2015deceptive}; \citet{correll2020truncating} \\
    Lines for trends, bars for separate values & \citet{zacks1999bars} \\
    One axis per measure & \citet{isenberg2011study}; \citet{haroz2016connected} \\
    The interval of the comparison, not two intervals & \citet{cumming2005inference}; \citet{belia2005researchers} \\
    Runs drawn under an estimate & \citet{hofman2020visualizing}; \citet{zhang2023illusion} \\
    Points for a handful of runs & \citet{weissgerber2015beyond,weissgerber2019reveal} \\
    Distribution display by what is judged & \citet{lotteraner2026investigating}; \citet{newburger2023fitting} \\
    Arrows for flow; no crossings & \citet{heiser2006arrows}; \citet{purchase1997aesthetic} \\
    Unbroken regions; a title on every picture & \citet{saket2014node}; \citet{haroz2015isotype} \\
    A linear scale, not a radar chart & \citet{diehl2010radial}; \citet{fuchs2017glyphs} \\
    Enlarged regions shown together & \citet{plumlee2006zooming}; \citet{lekschas2020pattern} \\
    Few category colors in a map of points & \citet{tseng2023categorical} \\
    \pvsection{Supported by methodological arguments} \\
    Ratings as answers, not only a mean & \citet{liddell2018analyzing} \\
    Precision--recall when positives are rare & \citet{saito2015precision} \\
    Reliability diagrams with the cases per bin & \citet{guo2017calibration} \\
    Confusion rows as shares of each class & \citet{gortler2022neo} \\
    Weights shown without a claim of cause & \citet{jain2019attention} \\
    No reading of distances in t-SNE maps & \citet{wattenberg2016tsne} \\
    \pvsection{Supported by our checks (Section~\ref{sec:checks})} \\
    Colormaps ordered in lightness & Table~\ref{tab:colormaps}, Figure~\ref{fig:lightness} \\
    Table fills at a fixed lightness & Table~\ref{tab:fills}, Figure~\ref{fig:fills} \\
    Diagram blocks named by label, not by fill & Figure~\ref{fig:palettes} \\
    Spectrogram levels in even lightness steps & Table~\ref{tab:colormaps} \\
    \pvsection{Our judgment, untested} \\
    Best value in bold with a fill & by analogy with \citet{strobelt2016guidelines} \\
    A key under each table & by inference from \citet{correll2014error} \\
    Better values darker in every column & against \citet{schloss2019mapping} \\
    The accent for the paper's own method & by analogy with \citet{wolfe2004attributes} \\
    Claim first in each caption & by inference from \citet{kong2018frames} \\
    Diverging stacked bars; ties set aside & after \citet{heiberger2014design} \\
    Fitted and held-out points told apart & after \citet{openai2023gpt4} \\
    Overlay opacity rising with the weight & by inference from \citet{schloss2019mapping} \\
    Spectrogram floor of $-80$~dB & a tool default \citep{librosa2026source} \\
    Sections with a gap, a name and a band & after \citet{bae2014reinforcing} \\
    Rounded links, open arrowheads, no borders & none \\
    Thresholds: six groups, eight lines, ten runs, six classes & none \\
    \bottomrule
  \end{pvtable}
\end{table}

\section{Drawing the exhibit}\label{sec:design}

\subsection{One accent color}

Every color in pubviz is derived from one accent color, in OKLab, a color
space built so that equal distances look about equally different
\citep{ottosson2020oklab}. The accent gives a strong color for the paper's
own method in figures; three fills for the best, second and third values of
a table, at fixed lightness levels; a band for section rows; a nine-step ramp
for shading cells and heatmaps, with lightness in one direction; and six
light fills for kinds of diagram blocks. Green and red for changes that help
or hurt are set apart from the accent. The LaTeX package and the Python
library compute the same colors from the same rules, and a test checks that
they agree for every role.

Fixed lightness levels are a choice the studies do not make for us.
Designer ramps grown from one seed color were read at least as well as
straight interpolations and were rated more pleasant
\citep[Sec.~4.1]{smart2020color}, and guidance on ramps puts the order in
lightness steps \citep[p.~29]{harrower2003colorbrewer}, but neither says how
a table's fills should be placed. Our check in Section~\ref{sec:checks} shows
what fixed levels buy: a fill made by mixing the accent with white, the usual
\code{blue!40} of LaTeX, changes its lightness with the hue, while a fill at a
fixed lightness looks the same for every accent. We set the best fill at a
lightness of 0.85, which in our check keeps the best and second fills apart
in grey for every preset accent while black text on them stays easy to read.

\subsection{Marks and the key}

Inside a table, the package finds the numbers, reads the direction of each
measure from an arrow where the measure is named (\code{Acc~\textbackslash
pvup}, \code{WER~\textbackslash pvdown}), ranks the values, and marks them:
bold and the first fill for the best value, the second fill for the next
(Table~\ref{tab:modes}). Ties share a mark. With \code{delta=absolute} or
\code{delta=relative}, a row marked as the baseline makes every other value
carry its signed change, colored by whether it helps. The
key under the table, \code{\textbackslash pvkey}, has one line for each kind
of mark the table actually shows, and each line opens with the mark itself,
so the caption is free to state the finding.

\begin{table}[tb]
  \centering\scriptsize
  \setlength{\tabcolsep}{2.4pt}
  \newcommand\demo[1]{%
    \begin{minipage}[t]{0.325\linewidth}\centering
      \textsf{mode=#1}\par\smallskip
      \pvsetup{mode=#1}%
      \begin{pvtable}[delta=absolute]{lccc}
        \toprule
        Method & Acc \pvup & WER \pvdown & RTF \pvdown \\
        \midrule
        \pvbaseline Base & 70.1 & 12.4 & 0.31 \\
        Prior A & 72.3 & 11.9 & 0.35 \\
        Prior B & 71.6 & 12.8 & 0.22 \\
        \pvours Ours & 74.0 & 10.2 & 0.28 \\
        \bottomrule
      \end{pvtable}
      \pvkey[baseline=Base]
    \end{minipage}}
  \caption{\textbf{The same table keeps its meaning in color, in grey and
    without color.} The package marks the best value in bold with a fill and
    the second with a lighter fill, adds each change against the baseline,
    and writes the key from what the table shows. Illustrative numbers.}
  \label{tab:modes}
  \demo{color}\hfill\demo{gray}\hfill\demo{plain}
\end{table}

These marks rest partly on analogy. Bold with a background fill is the pair
\citet[Sec.~8.1]{strobelt2016guidelines} recommend for highlighted words, and
when a display varies in both lightness and hue, viewers group it by
lightness first \citep[p.~1176]{healey2012attention}; we take this as one
reason to place the fills by lightness. Every color cue has a
second cue that needs no color: the best value is bold as well as filled, a
change carries its sign, and in \code{mode=plain} the second value is
underlined. Printing the sign matters because the meaning of red is learned:
Chinese stockbrokers, for whom red marks a rising price, did better after a
red cue, while students did worse \citep[Results]{zhang2014experience}.

\subsection{Grey and plain print}

A venue that prints in black and white gets \code{mode=gray}, in which every
color becomes the grey of the same lightness, and \code{mode=plain}, which
drops the fills for bold and underline. Diverging schemes do not survive a
black-and-white copy \citep[p.~30]{harrower2003colorbrewer}, and no method of
turning color images to grey was best for every image
\citep[Sec.~4.2]{cadik2008perceptual}; pubviz therefore designs the grey
version rather than leaving it to a printer, and orders every scale it uses
in lightness. APT already chose different designs for color and monochrome
output \citep[Sec.~9]{mackinlay1986apt}; pubviz applies one switch to
tables, figures and diagrams together.

\subsection{Uncertainty}

Three figure kinds follow the last chart of the tree. \code{pv.deltas} draws
the difference from a baseline with the interval of that difference, and
\code{pv.mean\_ci} gives the interval of per-seed changes when runs are
paired. \code{pv.means} draws the level of each method as a point and an
interval, or as a gradient, over the runs drawn faintly underneath.
\code{pv.dots} draws every run of a handful as a dot, moving dots that would
cover each other, and \code{pv.distributions} draws many runs as a box plot,
a strip, a violin or a histogram, chosen by what the reader should judge; it
warns when a method has too few runs to summarize. None of the three uses a
bar for an estimate.

\subsection{Groups within a figure}

Results often have two levels: the settings of one model, or the variants of
one ablated part. pubviz groups them into sections, as a table does with its
section rows, with three cues that agree: space before each section, its name
beside or under its members, and a light grey band behind it
(Figure~\ref{fig:sections}). An enclosing region is among the strongest
grouping cues: when it conflicted with spacing, 21 of 22 observers grouped by
the region \citep[p.~352]{montoro2015competicion}, and when people rebuilt
two-level structures from web-page layouts, enclosing boxes helped most and
cues that agreed beat cues that competed
\citep[Table~5, Sec.~4.4]{bae2014reinforcing}. pubviz therefore gives every
section a band by default rather than every other one. Grouping has a price:
it speeds comparisons within a group mainly by slowing those across groups
\citep[Exp.~3]{nothelfer2017redundant}, and gaps between groups make the
overall extreme less obvious \citep[Sec.~5.5]{fygenson2024arrangement}, so
sections suit a claim about the groups. No study has tested bands in charts.
The band is a neutral grey, so it is not read as a mark of ours, and it is
lighter than the tints publishers ask for in print \citep{ams2017graphics}; a
printer may drop it, and the space and the names still carry the groups.

\begin{figure}[tbp]
  \centering
  \begin{subfigure}{\linewidth}\centering
    \fig{catalog/bars-sections}\caption{\code{pv.bars} with sections of benchmarks}
  \end{subfigure}\\[4pt]
  \begin{subfigure}{0.485\linewidth}\centering
    \fig{catalog/means-sections}\caption{\code{pv.means} with a model and its settings}
  \end{subfigure}
  \caption{\textbf{Space, a name and a band mark each section, and agree.}
    The second figures of the repository's \code{bars.py} and
    \code{means.py}; illustrative numbers. In (b) the rows are named as
    (model, setting) pairs, so a setting's name can repeat across models.}
  \label{fig:sections}
\end{figure}

\subsection{Pictures, sound and weights}

The exhibits of a model at work keep the rules of the rest. Every scale
orders its values in lightness: the heatmap ramp for confusion matrices and
attention, and for spectrograms and the difference maps under images a grey
ramp from white to near black, which prints in grey unchanged. Ours keeps the
accent in classifier curves, and the pitch contour on a spectrogram takes it
too, since the spectrogram itself has no hue. A saliency map is laid over the
image in grey, so the only color on the image is the map's. Signed values use
the two hues for changes that help or hurt, and keep a second cue that needs
no color: a signed word shaded with \code{\textbackslash pvshade*} is also
underlined when its weight is negative. Axis titles such as ``True class'' or
``Recall'' follow the paper's language.

\subsection{Sizes from the paper's own layout}

A figure built larger than its column and shrunk to fit prints its labels
smaller than the text around them. pubviz avoids this by building each
figure at its printed width. Each compile of the paper writes a layout file,
\code{<jobname>.pubviz.json}, with the column width, the text width, the
body font size, the font families, the colors and the words of the key.
\code{pv.use} reads it, sizes the figure text a little below the body text,
and colors the figures from the same accent as the tables; \code{pv.save}
warns when a figure's width changed after it was made, when it has an axes
title that belongs in the caption, when text falls below five points, or
when texts overlap or run off the figure. The figure functions widen their
axes so that the values they print stay inside.
Because the layout comes from the paper's own compilation, the figures follow
any template without a table of venue sizes.

\subsection{Languages}

The words of a key, and the few words a figure writes itself, follow the
paper's language. English is the default, Traditional and Simplified Chinese,
Japanese and Korean are built in, and any other language is a list of words
given to \code{\textbackslash pvdefinelanguage}; the punctuation, and whether
words are separated by spaces, are part of that list. The layout file carries
the words to the figures, so a language defined in LaTeX reaches them too.
Chinese, Japanese and Korean characters in a figure fall back to an installed
font for that language and then to those of the others.

\subsection{Diagrams}

With \code{\textbackslash usepackage[diagrams]\{pubviz\}} the package adds
TikZ styles for method diagrams: blocks, operations, flow arrows and groups.
By default only the part a paper adds takes the accent
(Figure~\ref{fig:transformer}); with \code{pv kinds} each kind of block gets a
light fill of its own. Those fills are too close to tell apart under color
vision deficiency and in grey (Section~\ref{sec:checks}), so the label inside
each block, not its fill, names it. The diagram environment measures its
width and warns when it is wider than the line, since scaling it down would
shrink its text as well.

For stages and flows, pubviz also has a lighter style
(Figure~\ref{fig:light-diagrams}): a stack of layers drawn as a small picture
with a title above it, thin grey links with open arrowheads, and a light
region behind the parts that belong together. Method diagrams are close to
universal in machine learning papers: 82\% of the long papers at one NLP
conference had a diagram that was not a result, and 64\% a system diagram
\citep[Sec.~4.1]{marshall2021number}. Little tests their style. Text labels
were read with fewer errors than pictographs standing in for them
\citep[Exp.~1]{haroz2015isotype}, so every picture keeps its title, and readers
of neural network diagrams disagree on whether a drawn count is meant
\citep[Table~3]{marshall2025evidence}, so the real count goes in the grey
note. Unbroken group regions did not slow readers tracing paths
\citep[Sec.~5]{saket2014node}, while fragmented ones with labels on top did
\citep[Sec.~IV.B]{jianu2014how}, and flowcharts were followed best left to
right \citep[abstract]{krohn1983flowcharts}. Rounded corners, open arrowheads
and blocks without borders are our choice; no study has tested them.

\begin{figure}[tb]
  \centering
  \transformer[pv kinds=false, ffn/.append style={pv new}]
  \caption{\textbf{Only the part a paper adds takes the accent.} A
    Transformer drawn with the pubviz styles, with the feed-forward block
    taken as the new part. With \code{pv kinds}, each kind of block would get
    a light fill of its own instead.}
  \label{fig:transformer}
\end{figure}

\begin{figure}[tbp]
  \centering
  \begin{subfigure}{\linewidth}\centering
    \looped\caption{Blocks that run several times with shared weights}
  \end{subfigure}\\[8pt]
  \begin{subfigure}{\linewidth}\centering
    \asr\caption{A pipeline from speech to text}
  \end{subfigure}\\[8pt]
  \begin{subfigure}{\linewidth}\centering
    \retrieval\caption{Retrieval before generation}
  \end{subfigure}
  \caption{\textbf{A lighter style draws stages and flows with few lines.}
    Stacks of layers as small pictures, thin grey links with open arrowheads,
    a title and a grey note above each stage, and a light region behind the
    parts that belong together. The part each example takes as new is in
    the accent; in grey and plain print it turns near black with a heavier
    line.}
  \label{fig:light-diagrams}
\end{figure}

\section{The exhibits}\label{sec:catalog}

Every exhibit the tree can reach has a pubviz call (Table~\ref{tab:calls}),
and every figure kind comes with a short script that draws the examples in
Figures~\ref{fig:catalog}--\ref{fig:catalog-media} from their data to the
LaTeX that includes them.

\begin{table}[tb]
  \centering\small
  \caption{\textbf{Every answer of the tree has a call that draws it.} The
    exhibits of pubviz and the chart of the decision tree that leads to
    each.}
  \label{tab:calls}
  \begin{pvtable}[ranks=0, stretch=1.05]{lll}
    \toprule
    Exhibit & Call & Chart \\
    \midrule
    Result table: ranks, changes, shading, key & \code{pvtable}, \code{\textbackslash pvkey} & 2 \\
    Method diagram & \code{pvdiagram}, \code{pv} styles & 1 \\
    Bars from zero & \code{pv.bars} & 3 \\
    Differences with their intervals & \code{pv.deltas} & 3, 4 \\
    Series over steps or time & \code{pv.curves} & 1 \\
    A measure against size, with a fitted law & \code{pv.scaling} & 1 \\
    Answers on a rating scale & \code{pv.ratings} & 1 \\
    Trade-off with its frontier & \code{pv.tradeoff} & 1 \\
    Many items, one value each & \code{pv.ranked} & 3 \\
    Values over two ordered factors & \code{pv.heatmap} & 3 \\
    Every run of a handful & \code{pv.dots} & 4 \\
    Level with the uncertainty of the mean & \code{pv.means} & 4 \\
    Spread of many runs & \code{pv.distributions} & 4 \\
    Confusion matrix & \code{pv.confusion} & 5 \\
    Errors at every threshold & \code{pv.roc}, \code{pv.precision\_recall}, \code{pv.det} & 5 \\
    Reliability diagram & \code{pv.reliability} & 5 \\
    Map of embedded points & \code{pv.embedding} & 5 \\
    Image grid with an enlarged region & \code{pv.images} & 5 \\
    Spectrogram and waveform & \code{pv.spectrogram}, \code{pv.waveform} & 5 \\
    Weights between two sequences & \code{pv.attention} & 5 \\
    Weights over an image & \code{pv.overlay} & 5 \\
    Weights over words & \code{\textbackslash pvshade}, \code{pv.pvshade} & 5 \\
    \bottomrule
  \end{pvtable}
\end{table}

\begin{figure}[tbp]
  \centering
  \begin{subfigure}{\linewidth}\centering
    \fig{catalog/bars}\caption{\code{pv.bars}}
  \end{subfigure}\\[4pt]
  \begin{subfigure}{0.485\linewidth}\centering
    \fig{catalog/deltas}\caption{\code{pv.deltas}}
  \end{subfigure}\hfill
  \begin{subfigure}{0.485\linewidth}\centering
    \fig{catalog/curves}\caption{\code{pv.curves}}
  \end{subfigure}\\[4pt]
  \begin{subfigure}{0.485\linewidth}\centering
    \fig{catalog/tradeoff}\caption{\code{pv.tradeoff}}
  \end{subfigure}\hfill
  \begin{subfigure}{0.485\linewidth}\centering
    \fig{catalog/heatmap}\caption{\code{pv.heatmap}}
  \end{subfigure}
  \caption{\textbf{Each figure kind marks the paper's own method in the
    accent and every compared method in grey.} The examples of the
    repository for results measured once and for series, drawn at the width
    shown here; illustrative numbers.}
  \label{fig:catalog}
\end{figure}

\begin{figure}[tbp]
  \centering
  \begin{subfigure}{0.485\linewidth}\centering
    \fig{catalog/dots}\caption{\code{pv.dots}}
  \end{subfigure}\hfill
  \begin{subfigure}{0.485\linewidth}\centering
    \fig{catalog/means}\caption{\code{pv.means}}
  \end{subfigure}\\[4pt]
  \begin{subfigure}{0.485\linewidth}\centering
    \fig{catalog/distributions}\caption{\code{pv.distributions}, box}
  \end{subfigure}\hfill
  \begin{subfigure}{0.485\linewidth}\centering
    \fig{catalog/distributions-strip}\caption{\code{pv.distributions}, strip}
  \end{subfigure}
  \caption{\textbf{The figures for repeated runs show the runs themselves.}
    The examples of the repository for estimates over seeds and for
    distributions; illustrative numbers. The box plot in (c) hides the
    second, slow mode of the baseline that the strip in (d) shows.}
  \label{fig:catalog-runs}
\end{figure}

\begin{figure}[tbp]
  \centering
  \begin{subfigure}{0.485\linewidth}\centering
    \fig{catalog/ranked}\caption{\code{pv.ranked}}
  \end{subfigure}\hfill
  \begin{subfigure}{0.485\linewidth}\centering
    \fig{catalog/scaling}\caption{\code{pv.scaling}}
  \end{subfigure}\\[4pt]
  \begin{subfigure}{0.485\linewidth}\centering
    \fig{catalog/ratings}\caption{\code{pv.ratings}}
  \end{subfigure}\hfill
  \begin{subfigure}{0.485\linewidth}\centering
    \fig{catalog/ratings-preference}\caption{\code{pv.ratings}, ties aside}
  \end{subfigure}
  \caption{\textbf{Many items become a sorted dot chart, scale a fitted
    law on log axes, and ratings the share of each answer.} The examples of
    the repository for results over many items, over model size and from
    human judgments; illustrative numbers. In (b) the hollow point was held
    out of the fit.}
  \label{fig:catalog-scale}
\end{figure}

\begin{figure}[tbp]
  \centering
  \begin{subfigure}{0.485\linewidth}\centering
    \fig{catalog/confusion}\caption{\code{pv.confusion}}
  \end{subfigure}\hfill
  \begin{subfigure}{0.485\linewidth}\centering
    \fig{catalog/roc}\caption{\code{pv.roc}}
  \end{subfigure}\\[4pt]
  \begin{subfigure}{0.485\linewidth}\centering
    \fig{catalog/precision_recall}\caption{\code{pv.precision\_recall}}
  \end{subfigure}\hfill
  \begin{subfigure}{0.485\linewidth}\centering
    \fig{catalog/det}\caption{\code{pv.det}}
  \end{subfigure}
  \caption{\textbf{The same scores look strong as an ROC curve and weak as
    a precision--recall curve when positives are rare.} The examples of the
    repository for a model's errors; (b) and (c) draw the same made-up scores
    with 5\% positives.}
  \label{fig:catalog-model}
\end{figure}

\begin{figure}[tbp]
  \centering
  \begin{subfigure}{0.485\linewidth}\centering
    \fig{catalog/reliability}\caption{\code{pv.reliability}}
  \end{subfigure}\hfill
  \begin{subfigure}{0.485\linewidth}\centering
    \fig{catalog/embedding}\caption{\code{pv.embedding}}
  \end{subfigure}\\[4pt]
  \begin{subfigure}{0.485\linewidth}\centering
    \fig{catalog/attention}\caption{\code{pv.attention}}
  \end{subfigure}\hfill
  \begin{subfigure}{0.485\linewidth}\centering
    \fig{catalog/overlay}\caption{\code{pv.overlay}}
  \end{subfigure}\\[6pt]
  \begin{subfigure}{\linewidth}\centering\small
    \pvshade{0.08}{the} \pvshade{1.00}{cat} \pvshade{0.48}{sat} \pvshade{0.13}{on}
    \pvshade{0.05}{the} \pvshade{0.73}{mat} \pvshade{0.03}{.}\qquad
    \pvshade*{0.03}{The} \pvshade*{0.11}{film} \pvshade*{0.03}{was} \pvshade*{-0.58}{not}
    \pvshade*{-0.84}{bad} \pvshade*{0.05}{at} \pvshade*{0.08}{all}\pvshade*{0.01}{,}
    \pvshade*{0.21}{acting} \pvshade*{1.00}{superb}
    \caption{\code{\textbackslash pvshade} and \code{\textbackslash pvshade*}, set in this report's own text}
  \end{subfigure}
  \caption{\textbf{Probabilities, places and weights each keep one scale and
    no claim of cause.} The examples of the repository for whether a model's
    probabilities hold, where it places its inputs and what its outputs rest
    on; the points, weights and maps are made up.}
  \label{fig:catalog-weights}
\end{figure}

\begin{figure}[tbp]
  \centering
  \begin{subfigure}{0.485\linewidth}\centering
    \fig{catalog/images}\caption{\code{pv.images}}
  \end{subfigure}\hfill
  \begin{subfigure}{0.485\linewidth}\centering
    \fig{catalog/spectrogram}\caption{\code{pv.spectrogram}, \code{pv.waveform}}
  \end{subfigure}
  \caption{\textbf{Outputs are compared region by region and level by
    level.} The examples of the repository for a model's outputs; the images
    and the sound are made up. In (a) the boxed region is enlarged under each
    image, and the bottom row is each output's difference from the
    reference.}
  \label{fig:catalog-media}
\end{figure}

\section{Our checks}\label{sec:checks}

The studies leave three color choices open, and we checked them by
computation. All measures use OKLab \citep{ottosson2020oklab}, with lightness
$L$ in points from 0 (black) to 100 (white); text contrast is the WCAG ratio,
from 1 to 21, where 7 is the threshold for enhanced readability
\citep{w3c2008wcag}. Color vision deficiency is simulated with the matrices
of \citet{machado2009cvd} at full severity, which we apply to linear RGB
because the authors do not say which RGB they assume. A script in the
repository computes every value.

\paragraph{Colormaps.} A colormap whose lightness turns back reorders its
values in a grey print. Table~\ref{tab:colormaps} measures eight maps: jet and
turbo turn back by more than 50 points, while viridis, magma, Greys and the
three maps of pubviz never do (Figure~\ref{fig:lightness}). No map is best on
every measure. The table shading of pubviz spans only about 20 points so that
black text stays readable on every cell, and its figure map spans twice that
for heatmaps, where the value is printed in each cell. Its level map, for
spectrograms and difference maps, gives up black text for a span of about 80
points, close to that of Greys, while its steps stay far more even than those
of Greys.

\begin{table}[tb]
  \centering\small
  \caption{\textbf{Only the rainbow maps turn back in lightness, and no map is
    best on every measure.} Reversal: the largest move of lightness against
    the map's direction; span: the range of lightness ($L$ points);
    unevenness: variation of the color step between equal data steps; black
    text: the lowest WCAG contrast of black text on the map.}
  \label{tab:colormaps}
  \gen{colormaps}
\end{table}

\paragraph{Categorical palettes.} A palette that merges for a reader with a
color vision deficiency, or in grey, loses the identity of its series.
Figure~\ref{fig:palettes} gives the smallest difference between any two of
six colors. The first six colors of matplotlib nearly merge under simulated
protanopia; the Okabe--Ito colors keep apart under both deficiencies but
include two of the same lightness, which a grey print merges; and the light
kind fills of pubviz are close in every condition, which is why a
diagram's labels must carry the kinds.

\begin{figure}[tb]
  \centering
  \begin{minipage}[t]{0.485\linewidth}\centering
    \fig{lightness}
    \captionof{figure}{\textbf{Rainbow maps go dark, light and dark again;
      the others change lightness in one direction.} OKLab lightness along
      each colormap; colored lines: the three maps of pubviz.}
    \label{fig:lightness}
  \end{minipage}\hfill
  \begin{minipage}[t]{0.485\linewidth}\centering
    \fig{palettes}
    \captionof{figure}{\textbf{Common palettes merge under color vision
      deficiency or in grey.} Smallest OKLab difference between any two of
      six colors ($\times 100$); bold: the largest in each column.}
    \label{fig:palettes}
  \end{minipage}
\end{figure}

\paragraph{Table fills.} A fill made by mixing the accent with white changes
its lightness with the hue: across the seven accents of pubviz, the best
fill of a 40/15 mix varies by 9 points (Table~\ref{tab:fills}). Fills at a
fixed lightness do not vary, and they expose a trade-off between how far
apart the best and second fills stay in grey and how well black text reads
on the best fill (Figure~\ref{fig:fills}). Every mix we tried is beaten on
both counts by some fixed level. At 0.90 the grey gap shrinks to about five
points; at 0.85, the level pubviz now uses, it is about ten, with black
text still above a contrast of 12.

\begin{figure}[tb]
  \centering
  \begin{minipage}[t]{0.485\linewidth}\centering
    \vspace{0pt}\footnotesize\setlength{\tabcolsep}{3pt}
    \captionof{table}{\textbf{Fixed-lightness fills look the same for every
      accent.} Spread: range of the best fill's lightness over seven accents;
      grey gap (best against second fill) and contrast (black text on the
      best fill): the worst accent.}
    \label{tab:fills}
    \gen{fills}
  \end{minipage}\hfill
  \begin{minipage}[t]{0.485\linewidth}\centering
    \vspace{0pt}\fig{fills}
    \captionof{figure}{\textbf{Fixed levels beat every mix with white, and
      the level trades grey separation for text contrast.} Worst case over
      seven accents; the dashed line joins the designs no other design beats
      on both measures.}
    \label{fig:fills}
  \end{minipage}
\end{figure}

\paragraph{The software.} A test suite compiles the package with pdfLaTeX,
LuaLaTeX and XeLaTeX, checks that the LaTeX and Python sides derive the same
color for every role, checks each line of the key in several languages, and
draws every example figure, failing on any problem \code{pv.save} reports. The
venue presets of the Python library were measured from the layout files of
the venues' own classes where we could compile them.

\section{Limits}\label{sec:limits}

The evidence covers less than the tree. Most experiments used simple stimuli
with crowd workers or students, such as one ratio between two marks or
isolated color patches, and how far their effects reach expert readers of
real papers is untested. No study we found tested how a table should mark its
best value, a key under a table, a display of trade-offs, twin y axes for two
measures, or figures of three to ten runs, the usual number of seeds; these
answers in the tree are our judgment, and Table~\ref{tab:status} lists them.
The exhibits of a model at work rest on still less. For confusion matrices,
classifier curves, reliability diagrams, attention matrices and spectrograms
we found methodological arguments and tool conventions, not studies of how
readers take these displays in a paper; the studies of image comparison and
insets used interactive viewers or radiographs, and the counts of figure
types we cite cover no speech venue.
The thresholds in the tree are rules of thumb. Our checks are computations on
a color model, not studies with readers: they assume that OKLab lightness
predicts what a grey print keeps and that the deficiency matrices apply to
linear RGB. We chose the studies by topic, without a search protocol, and the
related work in Section~\ref{sec:related} covers the items we read. The words
of the key in Japanese and Korean have yet to be checked by native speakers.

\section{Conclusion}\label{sec:conclusion}

Research on graphical perception says a good deal about what an exhibit
should draw and much less about how a paper should mark and color it.
pubviz puts the first part into a decision tree, which reaches from result
tables to the exhibits of a model at work, and into the defaults of a LaTeX
package and a Python library, and marks the second part as judgment,
so a researcher can see which choices rest on readers and which on us. The
places where the tree rests on judgment, above all the marking of tables, are
where a study with readers would help most.

\bibliographystyle{plainnat}
\bibliography{refs}

\end{document}
